\section{Experimental Setup}
\label{sec:setup}

\paragraph{Environments.}
Each environment is an SCM over $K$ scalar variables with lagged
additive-noise mechanisms $z^i_{t+1} = \tanh\big(\sum_j W_{ij} z^j_t\big) + \epsilon^i_t$,
$\epsilon^i_t \sim \mathcal{N}(0, 0.1^2)$, on an Erd\H{o}s--R\'enyi graph with
edge probability $0.3$ and edge weights drawn uniformly from $\pm[0.5, 2]$.
Observations are $\x_t = \z_t + \mathcal{N}(0, 0.1^2)$. We use five graphs
with $K = 10$ (ER-10, 8 to 15 edges) and two with $K = 20$ (ER-20, 56 and 66
edges), generated from graph seeds $0$--$4$ and $0$--$1$. Training data is one
trajectory of $N = 10{,}000$ steps in which, at each step with probability
$0.5$, one random variable is set by $\doop(z^i = v)$, $v \sim U(-1.5, 1.5)$.
Each transition is paired with its counterfactual outcome under a different
random single-variable action and the same noise; whether a model uses these
pairs depends on the regime of Section~\ref{sec:arch}. For the data-efficiency
experiment we also train with $N = 2{,}000$ and $N = 500$ on the first three
ER-10 graphs.

\paragraph{Consequence task.}
Safety and goal are defined on variables downstream of the actions, so that a
forward model has to propagate an intervention through the graph. To avoid
hand-tuning, a fixed rule configures every environment: the safety variable is
the variable with the most ancestors in the true graph and the goal variable
the one with the second most; each threshold is the $0.7$-quantile of that
variable's realized three-step maximum under random single-variable actions
from states of the training process. An action is unsafe if the safety
variable reaches its threshold $\tau_s$ at any of the next $H = 3$ steps, and
reaches the goal if the goal variable exceeds $\tau_g$ at some step.
Each environment has 300 episodes with six candidate actions, all with values
in the training range, and outcomes are judged on the rollout of the true SCM
under one shared noise draw per episode. A quarter of the episodes contain no
safe candidate, so the correct response is to defer; every other episode
contains at least one safe candidate that reaches the goal, so the goal-rate
ceiling is $0.75$. Candidates whose realized outcome falls within $0.1$ of a
threshold are excluded, so that labels are decided by the action rather than
by the noise. The task comes in two versions: candidates that set one
variable, as in training, and candidates that set \emph{two} variables at
once, which never occurs in training. The episodes are identical for every
model. Table~\ref{tab:envs} in Appendix~\ref{app:envs} lists the resulting
configurations.

\paragraph{References and normalized violations.}
For every environment we evaluate two model-free planners: an \emph{oracle}
that simulates the noise-free true SCM from the noisy observation, and a
\emph{no-model} planner that predicts that nothing changes. Their violation
rates differ widely across environments (oracle $3.7\%$ to $16\%$, no-model
$27\%$ to $39\%$), so we report the \emph{normalized excess violation rate}
$(v - v_{\text{oracle}}) / (v_{\text{no-model}} - v_{\text{oracle}})$: $0$
for the oracle, $1$ for a planner that ignores the action. Raw rates are
given in Table~\ref{tab:envs}.

\paragraph{Rollout error.}
Independently of the planner, we measure the mean squared error of a model's
three-step prediction after a random one- or two-variable intervention,
against the noise-free true rollout from the true state, over 500 start
states drawn from the training process. Unlike one-step error on noisy
observations, which is dominated by noise for every model (Appendix~\ref{app:details}),
this isolates what the planner relies on.

\paragraph{Models.}
\emph{Vanilla VAE} uses one MLP on the concatenated latent and action;
\emph{Dreamer-style} uses a GRU cell; both share AC-LSCM's encoder and decoder.
\emph{Small Transformer} treats the $2K$ scalars of $(\x_t, a_t)$ as tokens
and predicts $\x_{t+1}$; it is a forward model plugged into the planner, not a
language-model agent. Absent action components are encoded as $0$ for all
baselines. Architectures and optimization settings are those of the first
version (Appendix~\ref{app:details}); none were tuned on the agent task. Each
baseline is evaluated as is (\emph{native}) and with the do-wrapper. Latent
models roll out in their latent space and decode every step.

\paragraph{Seeds and statistics.}
Each model is trained with three seeds per environment (two at $N < 10{,}000$),
408 training runs in total. The seed fixes the training trajectory and the
initialization, and the same seeds are used for every model, so comparisons
between models are paired. Tables report the mean over environments of the
per-environment seed means, with the standard deviation across environments,
which mostly reflects differences in difficulty between graphs. Paired
comparisons use per-(environment, seed) differences.

\section{Results}

\subsection{Main comparison}
\label{sec:main}

\begin{table}[t]
\centering
\scriptsize
\setlength{\tabcolsep}{4pt}
\begin{tabular}{@{}lccccc@{}}
\toprule
& \multicolumn{2}{c}{Rollout MSE $\downarrow$}
& \multicolumn{2}{c}{Normalized excess violations $\downarrow$}
& Goal rate $\uparrow$ \\
\cmidrule(lr){2-3}\cmidrule(lr){4-5}\cmidrule(l){6-6}
Planner's forward model & 1 variable & 2 variables & 1 variable & 2 variables & 1 variable \\
\inputrows{generated/v3_main.tex}
\bottomrule
\end{tabular}
\caption{Main comparison ($N = 10{,}000$; mean over graphs of per-graph means
over 3 seeds, $\pm$ standard deviation across graphs; best non-reference
entry in bold). ``1/2 variables'': interventions on one variable (as in
training) or two variables at once (never seen in training). Normalized
excess violations: $0$ = oracle, $1$ = a planner that predicts no effect.
Goal-rate ceiling $0.75$. ``+ do-wrapper'': the same trained model with the
correct action semantics imposed from outside (Section~\ref{sec:planning}).
The true-graph model is a reference, not a competitor.}
\label{tab:main}
\end{table}

\begin{table}[t]
\centering
\scriptsize
\setlength{\tabcolsep}{4pt}
\begin{tabular}{@{}lcccc@{}}
\toprule
& \multicolumn{2}{c}{ER-10 (5 graphs)} & \multicolumn{2}{c}{ER-20 (2 graphs)} \\
\cmidrule(lr){2-3}\cmidrule(l){4-5}
Baseline & 1 variable & 2 variables & 1 variable & 2 variables \\
\midrule
\inputrows{generated/v3_paired_a.tex}
\midrule
& \multicolumn{2}{c}{ER-10, $N = 2{,}000$ (3 graphs)} & \multicolumn{2}{c}{ER-10, $N = 500$ (3 graphs)} \\
\cmidrule(lr){2-3}\cmidrule(l){4-5}
Baseline & 1 variable & 2 variables & 1 variable & 2 variables \\
\midrule
\inputrows{generated/v3_paired_b.tex}
\bottomrule
\end{tabular}
\caption{Paired comparison with the factored AC-LSCM: normalized excess
violations of the baseline minus those of the factored AC-LSCM trained on the
same data (positive: the factored model is safer), mean $\pm$ standard error
over (graph, seed) pairs, and in parentheses the number of pairs in which the
factored model is safer. In all $396$ corresponding comparisons of rollout
MSE, the factored model has the lower error.}
\label{tab:paired}
\end{table}

Table~\ref{tab:main} compares all planners at $N = 10{,}000$, and
Table~\ref{tab:paired} gives the paired comparison of the factored AC-LSCM
with each baseline.

\emph{The factored AC-LSCM is the most accurate forward model.} It has the
lowest three-step rollout error on both graph sizes and both action types,
and it has lower rollout error than the corresponding baseline in every one of
the $396$ paired comparisons (Table~\ref{tab:paired}). It recovers the true
graph exactly, with no missing or spurious edge, in all 15 runs on ER-10, and
misses 2 to 4 of 56 to 66 edges on ER-20, again with no spurious edge.

\emph{It is the safest planner, by margins that grow with the difficulty of
the task.} On ER-10 with single-variable actions, its normalized excess
violation rate is $0.25$, against $0.31$ for the best baseline (Vanilla VAE
with the do-wrapper); this difference is within the noise of the paired
comparison (8 of 15 pairs). With two-variable actions, which no model saw in
training, the factored model improves to $0.18$, while the baselines range
from $0.28$ to $0.55$, and it is safer than every baseline configuration in at
least 13 of 15 pairs. On ER-20 its excess violations are
$0.41$ against $0.53$ to $0.81$ for the baselines. In raw terms its violation rate is lower than that of every
baseline configuration, native or with the do-wrapper, in six of the seven
environments (Table~\ref{tab:envs}); other AC-LSCM variants are sometimes
lower still in individual environments. These gains do not come from deferring more: its goal
rate ($0.73$ on ER-10 against a ceiling of $0.75$) is the highest or within
$0.01$ of the highest.

\emph{Imposing the correct action semantics helps every baseline, but does not
close the gap.} The do-wrapper lowers the excess violations of Dreamer-style
from $0.49$ to $0.38$ and of the Transformer from $0.45$ to $0.36$ on ER-10,
and on ER-20 those of every baseline by $0.15$ to $0.25$. A large part of the
baselines' deficit is therefore in how they represent actions, and it can be
removed by any agent designer who knows which variables an action sets.
What remains after the do-wrapper is a difference in the learned dynamics,
and there the factorization is what helps.

\emph{The learned graph is nearly as good as the true one.} With the true
graph fixed, the factored model reaches $0.24$ on ER-10 (learned: $0.25$) and
$0.34$ on ER-20 (learned: $0.41$); the ER-20 gap comes from the few missed
edges.

\subsection{Why the latent AC-LSCM fails}
\label{sec:latent}

The same architecture with a learned encoder and decoder is the worst planner
on both graph sizes (excess violations $0.68$ on ER-10, $0.90$ on ER-20),
although on ER-10 its one-step prediction error is close to the baselines'. Two
observations locate the problem. First, its rollout error drops from $0.032$
to $0.012$ on ER-10 and from $0.157$ to $0.055$ on ER-20 when the do-wrapper
replaces its own action handling, and its excess violations from $0.68$ to
$0.34$. The model applies $\doop(z^i = v)$ to latent coordinate $i$, and
nothing in training makes that coordinate the observed variable $i$; what it
learns is a mechanism that reproduces the observed effect of the logged
actions through the decoder, which is enough for one-step prediction of
training-like transitions but not for multi-step simulation of new
interventions. Second, the graph it learns is a graph over these misaligned
coordinates: on average $7.9$ wrong edges on ER-10 and $62$ on ER-20, where the
factored model has $0$ and $3.3$. Even with the do-wrapper, its dynamics are
less accurate than the factored model's ($0.012$ versus $0.009$ on ER-10,
$0.055$ versus $0.032$ on ER-20). Supervising it with true counterfactual
latents, as in the first version, anchors the coordinates partly and lowers
its excess violations to $0.52$ on ER-10, still far from the factored
model, and does not help on ER-20 (Table~\ref{tab:ablation}).

The lesson is that the benefit of causal structure depends on placing it on
variables for which the intervention semantics are right. When the agent's
actions and predicates refer to observed variables, those variables are the
natural factors; learning a latent space and hoping that it aligns with the
action targets loses the benefit and adds a failure mode.

\subsection{Data efficiency}
\label{sec:data}

\begin{table}[t]
\centering
\scriptsize
\setlength{\tabcolsep}{2.1pt}
\begin{tabular}{@{}lccccccccc@{}}
\toprule
& \multicolumn{3}{c}{$N = 500$} & \multicolumn{3}{c}{$N = 2{,}000$} & \multicolumn{3}{c}{$N = 10{,}000$} \\
\cmidrule(lr){2-4}\cmidrule(lr){5-7}\cmidrule(l){8-10}
Model & Rollout & Excess (1) & Excess (2) & Rollout & Excess (1) & Excess (2) & Rollout & Excess (1) & Excess (2) \\
\midrule
\inputrows{generated/v3_data.tex}
\bottomrule
\end{tabular}
\caption{Training-set size (first three ER-10 graphs; 2 seeds at $N < 10{,}000$,
3 at $N = 10{,}000$; native action handling). Rollout: three-step rollout MSE
after single-variable interventions. Excess (1)/(2): normalized excess
violations with one-/two-variable candidate actions. Best non-reference entry
per column in bold.}
\label{tab:data}
\end{table}

Table~\ref{tab:data} varies the amount of training data. The factored model's
rollout error at $N = 2{,}000$ ($0.007$) is lower than that of every baseline
at $N = 10{,}000$, and at $N = 500$ it is $2.5$ to $5$ times lower than the
baselines' at the same size. In planning, its advantage with single-variable
actions at $N = 500$ is small (excess $0.31$ against $0.32$ for the VAE),
but with two-variable actions it is large at every data size ($0.29$ against
at best $0.48$ at $N = 500$; $0.15$ against $0.41$ at $N = 2{,}000$). At
$N = 500$ the sparsity penalty has not pruned any edge by the end of
training, so the factored model is effectively dense there; its advantage at
that size comes from the per-variable factorization and the action
semantics, not from a learned graph. At $N = 2{,}000$ it recovers the exact
graph on all three graphs.

\subsection{Ablations}
\label{sec:ablations}

\begin{table}[t]
\centering
\scriptsize
\setlength{\tabcolsep}{3.5pt}
\begin{tabular}{@{}llcccccc@{}}
\toprule
Model & CF data & SHD & 1-step MSE & CF MSE & Rollout & Excess (1) & Excess (2) \\
\inputrows{generated/v3_ablation.tex}
\bottomrule
\end{tabular}
\caption{Ablations ($N = 10{,}000$, native action handling, mean over graphs).
CF data: counterfactual training pairs given to none of the models, the
observed counterfactual outcome given to every model, or (${}^\dagger$, first
version) noise-free counterfactual states given to AC-LSCM only. SHD:
structural Hamming distance to the true graph (true edges: 8 to 15 on ER-10,
56 and 66 on ER-20). One-step and counterfactual MSE are on noisy held-out
observations.}
\label{tab:ablation}
\end{table}

Table~\ref{tab:ablation} varies one component at a time.

\emph{The explicit do-operator adds little once the factors are right.}
Replacing the hard do-operator of the factored model by the soft one, in
which the intervened variable's mechanism receives the action value and must
learn to output it, changes excess violations from $0.25$ to $0.26$ and leaves
rollout error unchanged. The mechanisms learn the action semantics from data.
This agrees with the do-operator ablation of the first version, and it is the
mirror image of the do-wrapper result: what matters is that the intervention
lands on the right variable, not the particular mechanism that implements it.

\emph{Sparsity helps, more on larger graphs.} Without the sparsity penalty
the factored model keeps all edges; its excess violations rise from $0.25$ to
$0.30$ on ER-10 and, with two-variable actions on ER-20, from $0.40$ to
$0.53$.

\emph{Counterfactual data helps structured models only.} Giving every model
the observed counterfactual outcomes leaves the factored model unchanged on
ER-10 and lowers its excess violations on ER-20 from $0.41$ to $0.33$, below
the true-graph reference. It also helps the dense latent model, which
reaches $0.20$ on ER-10, the lowest single-variable entry of any model in our
study (it was not run on ER-20). It does not help the baselines, whose excess
violations stay level or rise (Vanilla VAE $0.33 \to 0.33$, Dreamer-style
$0.49 \to 0.57$, Transformer $0.45 \to 0.43$ on ER-10). One explanation is
that a model with an abduction step can use a counterfactual pair as
information about the noise-free mechanism, whereas for a model without one
it is another transition from the same state with a different action and an
outcome that shares the factual noise.

\subsection{Distance to the oracle}

No planner comes close to the oracle. The best factored configurations
retain about a quarter of the oracle-to-no-model gap on ER-10 and a third on
ER-20, and so does the true-graph model. The remaining errors are therefore
not structural. They come from fitting the mechanisms from noisy
observations, and from starting the rollout at the noisy observation rather
than the true state (the oracle starts there too, but with noise-free
mechanisms). Planning under such errors would benefit from uncertainty
estimates and conservative safety margins, which our planner does not use.
